\section{A proper quadratic quotient theorem for the ternary/four norm cone}
Let $\lambda=\lambda_3$, $a=\lambda-1$, $b=1/\lambda$, and
$B=1+az+bz^2$, so $P_3=(1-\lambda z)B$.
For $E=[3]^r[4]^s$, use ACTUAL polynomial division
\[
                   E=BJ+R_0+R_1z.                         \tag{1}
\]
This is the proper polynomial quotient, not a zero-truncated inverse.
General finite actual mixed3/4 C, full A2 and original one-descent
coefficient unimodality remain OPEN.

\begin{theorem}
For every $r\ge32$, $s\ge0$ such that
\[
                         b^{r+1}4^s\le1/64,              \tag{2}
\]
$J$ is strictly positive on its entire coefficient interval and is
globally ordinary log-concave. Its two remainders satisfy
$|R_0|,|R_1|\le1/4$. In particular this holds for EVERY
\[
                         s\ge16,\qquad r\ge3s-4.         \tag{3}
\]
Both counts are unbounded. No new finite entrance is needed.
\end{theorem}

\subsection{Forward locality replaces reverse endpoint repairs}
The exact identity $[4]=2-P_3$ gives the COMPLETE quotient update
\[
 J_{r,s+1}=[4]J_{r,s}-(1-\lambda z)(R_0+R_1z),\qquad
 (R_0,R_1)_{r,s+1}=2(R_0,R_1)_{r,s}.                      \tag{4}
\]
The correction in (4) is confined to ranks ZERO, ONE, TWO.
Consequently every new ordinary minor with centre at least FOUR is
exactly the corresponding uniform-convolution minor. Only centres
one, two and three, and positivity at ranks zero through two, need
a new payment. No signed correction is discarded.

The previously proved ternary remainder map and its dual norm give
\[
 bR_0^2-aR_0R_1+R_1^2=b^{r+1}4^s.                       \tag{5}
\]
The initial remainder for $E=1$ is $(1,0)$, of norm $b$;
each ternary factor multiplies the norm by $b$ and each four factor
by FOUR. The generic complete two-natural-count induction is checked
again in the new kernel, and the earlier exact ternary map is reused.
The earlier coercivity theorem gives
\[
 (R_0^2+R_1^2)/4\le bR_0^2-aR_0R_1+R_1^2.
\]
Thus (2) puts BOTH remainders in the quarter box, with no phase
assumption. We retain the exact remainders rather than replacing
them by their norm in the polynomial update.

\subsection{Three robust head minors at all real counts}
Here is a field-independent head payment. Assume merely
$0\le a,b\le1$, $|R_0|,|R_1|\le1/4$, and $w=r+s\ge32$.
The first CI coefficients have the exact forms and bounds
\[
 E_1=w,\quad E_2=w(w+1)/2,
\]
\[
 E_3=w(w+1)(w+2)/6-r\in[U_3-w,U_3],\quad
 E_4=w(w+1)(w+2)(w+3)/24-rw-s\le U_4,
\]
where $U_3=w(w+1)(w+2)/6$, $U_4=w(w+1)(w+2)(w+3)/24$.
Coefficient comparison in (1) gives
\[
 J_0=1-R_0,\quad J_1=w-aJ_0-R_1,\quad
 J_i=E_i-aJ_{i-1}-bJ_{i-2}\ (2\le i\le4).
\]
Define
\[
 L_1=w-3/2,\quad L_2=E_2-w-3/2,\quad
 L_3=U_3-E_2-2w-1/4.
\]
The complete successive bounds are
\[
 3/4\le J_0\le5/4,\quad L_1\le J_1\le w+1/4,
 \quad L_2\le J_2\le E_2,
 \quad L_3\le J_3\le U_3,\quad J_4\le U_4.                \tag{6}
\]
For each step, a nonnegative mixture $ax+by$ lies between zero and
$x+y$. Positivity of each preceding coefficient follows from the
already paid lower bound; in particular $J_0,J_1,J_2>0$.
The THREE complete scalar budgets are
\[
 L_1^2\ge(5/4)E_2,\quad
 L_2^2\ge(w+1/4)U_3,\quad
 L_3^2\ge E_2U_4.                                      \tag{7}
\]
Every inequality in (7), and $L_i\ge0$, is a polynomial with
nonnegative rational coefficients after $w=32+x$, $x\ge0$.
This whole REAL-count statement is kernel checked; no bounded integer
sample is its proof. Combining (6)--(7) proves ALL three head minors
$J_1^2-J_0J_2$, $J_2^2-J_1J_3$, $J_3^2-J_2J_4\ge0$.
The kernel includes the full coefficient-bound assembly, not just
the three sufficient polynomials.

\subsection{A complete induction throughout the norm cone}
Fix any $r\ge32$. At $s=0$ use the EXISTING pure-ternary proper
quotient theorem, proved for every $r\ge18$ in
\texttt{tribonacci\_rich\_reverse\_quotient\_proof.tex}.
In particular this is a previously proved unbounded pure entrance,
not an assumed new diagonal.

Fix a target $s$ obeying (2). Every earlier four count has SMALLER
norm, since it is $b^{r+1}4^j$ with $j\le s$. Suppose $J_{r,j}$
is positive globally LC. Uniform convolution by $[4]$ is positive
globally LC. By (4), all coefficients at ranks at least three are
unchanged; all minor centres at least four are unchanged as well.
The new remainder is in the quarter box by its target norm, so
(6)--(7) pay the new ranks zero, one, two and all THREE missing
minor centres. This proves positive interval support and global LC
of $J_{r,j+1}$. The entire increasing-norm target induction is
kernel checked, including the NEW-norm guard used in this step.
It proves every target (2). No primitive LC is inferred from ordinary
LC of $J$.

For the explicit cone (3), the exact reciprocal-root bounds
$1/2\le b\le6/11$ give $4b^3\le2/3$. If $r+4\ge3s$, then
\[
 (b^{r+1}4^s)b^3=b^{r+4}4^s\le(4b^3)^s.
\]
Since $b^3\ge1/8$, for every $s\ge16$,
\[
 b^{r+1}4^s\le8(2/3)^s\le8(2/3)^{16}<1/64.
\]
Also $r\ge3s-4\ge44>32$. This proves the entire explicit cone,
with all natural exponents and threshold arithmetic kernel checked.

\paragraph{Exact scope.}
Nine new clean scoped reports cover the real-count head polynomials,
full three-head-minor assembly, unit-mixture bounds, proper four-factor
correction, norm scale, whole two-count norm, whole target induction,
generic interval-minor payment and the entire explicit norm cone.
The earlier pure-ternary quotient/coercivity/root certificates are
reused without recounting. Polynomial division and LC-convolution
identifications are written proofs. There is NO new finite proof
premise; extra exact tests are QA only. This theorem is not an actual
finite-arm derivative theorem. Its finite-count source and both
remainder coefficients must still be retained in a subsequent actual
ordinary-tail or C bridge.
